%% paper80_jacobian_vermutung_de.tex
%% Die Jacobian-Vermutung (Deutsch)
%% Autor: Michael Fuhrmann
%% Specialist Mathematics Project, Build 166
%% Datum: 2026-03-12

\documentclass[12pt,a4paper]{amsart}
\usepackage{amsmath,amssymb,amsthm,mathtools,enumitem,booktabs,array,hyperref}
\usepackage[utf8]{inputenc}
\usepackage[T1]{fontenc}
\usepackage[ngerman]{babel}
\usepackage{geometry}
\geometry{margin=2.5cm}

%% Theorem-Umgebungen (Deutsch)
\newtheorem{theorem}{Satz}[section]
\newtheorem{satz}[theorem]{Satz}
\newtheorem{lemma}[theorem]{Lemma}
\newtheorem{korollar}[theorem]{Korollar}
\newtheorem{corollary}[theorem]{Korollar}
\newtheorem{proposition}[theorem]{Proposition}
\newtheorem{vermutung}[theorem]{Vermutung}
\newtheorem{conjecture}[theorem]{Vermutung}
\theoremstyle{definition}
\newtheorem{definition}[theorem]{Definition}
\newtheorem{definition_de}[theorem]{Definition}
\newtheorem{beispiel}[theorem]{Beispiel}
\newtheorem{example}[theorem]{Beispiel}
\theoremstyle{remark}
\newtheorem{bemerkung}[theorem]{Bemerkung}
\newtheorem{remark}[theorem]{Bemerkung}

\hypersetup{
  colorlinks=true,
  linkcolor=blue,
  citecolor=blue,
  urlcolor=blue
}

\begin{document}

\title{Die Jacobian-Vermutung:\\
Polynomiale Automorphismen, Gradreduktion\\
und algebraische Äquivalenzen}

\author{Michael Fuhrmann}
\address{Specialist Mathematics Project, Build 166}
\email{specialist-maths@project.local}
\date{2026-03-12}

\begin{abstract}
Die Jacobian-Vermutung, gestellt von Keller im Jahr 1939, fragt, ob eine polynomiale
Abbildung $F \colon \mathbb{C}^n \to \mathbb{C}^n$ mit konstanter, von null verschiedener
Jacobi-Determinante $\det(JF) \in \mathbb{C}^*$ notwendigerweise ein polynomialer
Automorphismus ist. Trotz ihrer elementaren Formulierung bleibt die Vermutung für
$n \geq 2$ über $\mathbb{C}$ offen und gilt als eines der schwierigsten offenen Probleme
der affinen algebraischen Geometrie.

Wir präsentieren die zentralen Strukturresultate: die Bass--Connell--Wright-Reduktion
auf Abbildungen vom Grad 3, die Yagzhev-Symmetriereduktion sowie die bemerkenswerte
Äquivalenz mit der Dixmier-Vermutung über Automorphismen der Weyl-Algebra
$A_n(\mathbb{C})$, bewiesen von Tsuchimoto (2005) und Belov-Kanel--Kontsevich (2007).
Das reelle Analogon ist für $n \geq 2$ \emph{falsch} (Pinchuk 1994).
\end{abstract}

\maketitle

\tableofcontents

%% =============================================================================
\section{Einleitung und Geschichte}
%% =============================================================================

Im Jahr 1939 stellte O.~H.\ Keller~\cite{Keller1939} folgende Frage, heute bekannt
als \emph{Kellers Problem} oder die \textbf{Jacobian-Vermutung}:

\begin{vermutung}[Jacobian-Vermutung, Keller 1939]\label{verm:JV}
Sei $F = (F_1,\ldots,F_n) \colon \mathbb{C}^n \to \mathbb{C}^n$ eine polynomiale
Abbildung. Falls die Jacobi-Determinante
\[
  \det(JF)(x) = \det\!\left(\frac{\partial F_i}{\partial x_j}(x)\right)_{1\leq i,j\leq n}
\]
eine von null verschiedene Konstante ist, dann ist $F$ ein polynomialer Automorphismus,
d.h.\ es existiert eine polynomiale Abbildung $G \colon \mathbb{C}^n \to \mathbb{C}^n$
mit $F \circ G = G \circ F = \mathrm{id}$.
\end{vermutung}

Die Bedingung $\det(JF) = c \neq 0$ ist \emph{notwendig} für die Invertierbarkeit:
Ist $F$ invertierbar mit polynomialem Inversen $G$, so gilt nach der Kettenregel
$JF \cdot JG = I$, also $\det(JF) \cdot \det(JG) = 1$. Da beide Polynome sind und
ihr Produkt $1$ ergibt, müssen beide Konstanten sein.
Die Vermutung fragt, ob diese notwendige Bedingung auch hinreichend ist.

\subsection{Kellers ursprünglicher Kontext}

Keller arbeitete über $\mathbb{Z}$; seine Formulierung fragte, ob eine ganzzahlige
polynomiale Abbildung mit Jacobi-Determinante $\pm 1$ über $\mathbb{Z}$ invertierbar
sein muss. Die komplexe Version wurde schnell zur Standardform.

\subsection{Übersicht dieses Artikels}

Wir behandeln folgende Themen:
\begin{enumerate}[label=(\roman*)]
  \item Den trivialen Fall $n = 1$ (Abschnitt~\ref{abschn:n1}).
  \item Das Gradreduktionstheorem von Bass--Connell--Wright (Abschnitt~\ref{abschn:BCW}).
  \item Die Yagzhev-Symmetriereduktion (Abschnitt~\ref{abschn:yagzhev}).
  \item Struktureigenschaften polynomialer Automorphismen (Abschnitt~\ref{abschn:struct}).
  \item Die reelle Jacobian-Vermutung und Pinchuk's Gegenbeispiel (Abschnitt~\ref{abschn:real}).
  \item Bekannte Teilresultate für $n = 2$ (Abschnitt~\ref{abschn:n2}).
  \item Äquivalenz mit der Dixmier-Vermutung (Abschnitt~\ref{abschn:dixmier}).
\end{enumerate}

%% =============================================================================
\section{Der Fall \texorpdfstring{$n=1$}{n=1}: Eine vollständige Lösung}
\label{abschn:n1}
%% =============================================================================

\begin{satz}[Jacobian-Vermutung für $n=1$]\label{satz:n1}
Sei $f \colon \mathbb{C} \to \mathbb{C}$ eine polynomiale Abbildung mit $f'(x) = c \neq 0$
für alle $x \in \mathbb{C}$. Dann ist $f$ ein polynomialer Automorphismus.
\end{satz}

\begin{proof}
Die Bedingung $f'(x) = c$ für alle $x$ bedeutet, dass $f'$ ein konstantes Polynom
ist. Integration ergibt: $f(x) = cx + d$ mit $c, d \in \mathbb{C}$, $c \neq 0$.
Das Inverse ist $g(x) = c^{-1}(x - d)$, ein Polynom. Also ist $f$ ein polynomialer
Automorphismus.
\end{proof}

\begin{bemerkung}
Über Körpern der Charakteristik $p > 0$ versagt das Argument: Das Polynom
$f(x) = x + x^p$ hat $f'(x) = 1$ überall (da $(x^p)' = px^{p-1} = 0$), ist
aber als polynomiale Abbildung nicht invertierbar. Die Charakteristik-null-Voraussetzung
ist also wesentlich.
\end{bemerkung}

%% =============================================================================
\section{Polynomiale Automorphismen: Strukturtheorie}
\label{abschn:struct}
%% =============================================================================

\begin{definition_de}
Die \emph{Gruppe der polynomialen Automorphismen} $\mathrm{GA}_n(\mathbb{C})$ ist
die Gruppe aller bijektiven polynomialen Abbildungen $F \colon \mathbb{C}^n \to \mathbb{C}^n$
mit polynomialem Inversen.
\end{definition_de}

Zwei wichtige Untergruppen erzeugen einen Großteil der Struktur:

\begin{definition_de}
\begin{enumerate}[label=(\alph*)]
  \item Die \emph{affine Gruppe} $\mathrm{Aff}_n(\mathbb{C})$: Abbildungen
        $F(x) = Ax + b$ mit $A \in \mathrm{GL}_n(\mathbb{C})$, $b \in \mathbb{C}^n$.
  \item Die \emph{triangulare (Jonquières-)Gruppe} $\mathrm{BA}_n(\mathbb{C})$:
        Abbildungen der Form
        \[
          F(x_1,\ldots,x_n) = \bigl(a_1 x_1 + p_1,\;
           a_2 x_2 + p_2(x_1),\;\ldots,\;
           a_n x_n + p_n(x_1,\ldots,x_{n-1})\bigr)
        \]
        mit $a_i \in \mathbb{C}^*$ und $p_i \in \mathbb{C}[x_1,\ldots,x_{i-1}]$.
\end{enumerate}
\end{definition_de}

\begin{satz}[Jung--van der Kulk, 1942--1953]\label{satz:jung}
Für $n = 2$ ist jeder polynomiale Automorphismus von $\mathbb{C}^2$ eine endliche
Komposition affiner und triangularer Automorphismen:
\[
  \mathrm{GA}_2(\mathbb{C}) = \mathrm{Aff}_2(\mathbb{C}) *_{\mathrm{Aff}_2 \cap \mathrm{BA}_2} \mathrm{BA}_2(\mathbb{C}).
\]
\end{satz}

Diese amalgamierte Produktstruktur schlägt für $n \geq 3$ fehl
(Shestakov--Umirbaev~\cite{ShU04}: Nagata-Automorphismus ist nicht zahm).

%% =============================================================================
\section{Die Bass--Connell--Wright-Gradreduktion}
\label{abschn:BCW}
%% =============================================================================

Ein zentrales Vereinfachungsergebnis zeigt, dass die Jacobian-Vermutung für beliebige
Gradpolynome auf den Fall von kubisch-linearen Abbildungen reduziert werden kann.

\begin{definition_de}
Eine polynomiale Abbildung $F = (F_1,\ldots,F_n)$ heißt \emph{kubisch-linear}, falls
$F(x) = x + H(x)$, wobei jede Komponente $H_i$ ein homogenes Polynom dritten Grades ist.
\end{definition_de}

\begin{satz}[Bass--Connell--Wright, 1982]\label{satz:BCW}
Die Jacobian-Vermutung ist wahr für alle $n$ und alle polynomialen Abbildungen
genau dann, wenn sie für alle kubisch-linearen Abbildungen $F(x) = x + H(x)$ mit
$H$ homogen vom Grad $3$ und $\det(I + JH) = 1$ wahr ist.
\end{satz}

\begin{proof}[Beweisidee]
\textbf{Schritt 1}: Normierung. Durch Komposition mit $L^{-1} = (JF(0))^{-1}$
kann $F = \mathrm{id} + H$ mit verschwindendem Linearanteil von $H$ angenommen werden.

\textbf{Schritt 2}: Gradreduktion. Terme vom Grad $d > 3$ werden durch Einführung
neuer Variablen in kubische Terme umgewandelt (auf Kosten erhöhter Dimension).
Das algebraische Identitätsargument:
\[
  (JH)^2 = 0 \implies \det(I + JH) = 1,
\]
da nilpotente Matrizen nur den Eigenwert $0$ haben. Die Nilpotenz von $JH$ wird
unter der Gradreduktionseinbettung erhalten.
\end{proof}

\begin{korollar}[Nilpotente Reduktion]\label{kor:nilpotent}
Die Jacobian-Vermutung reduziert sich auf: Jede Abbildung $F = x + H(x)$ mit
kubisch-homogenem $H$ und nilpotentem $JH$ ist ein polynomialer Automorphismus.
\end{korollar}

%% =============================================================================
\section{Die Yagzhev-Symmetriereduktion}
\label{abschn:yagzhev}
%% =============================================================================

\begin{satz}[Yagzhev, 1980]\label{satz:yagzhev}
Die Jacobian-Vermutung ist äquivalent dazu: Jede Abbildung $F = x - H(x)$ mit
$H$ symmetrisch-kubisch-polynomial und $(JH)^2 = 0$ ist ein polynomialer Automorphismus.
\end{satz}

Die Symmetriebedingung bedeutet: die zugehörige multilineare Abbildung
$\hat{H} \colon (\mathbb{C}^n)^3 \to \mathbb{C}^n$ mit $H(x) = \hat{H}(x,x,x)$
ist symmetrisch in ihren drei Argumenten. Dies entspricht $H = \nabla \Phi$ für
ein Polynom $\Phi$ vom Grad $4$ (Gradientenabbildung).

%% =============================================================================
\section{Die reelle Jacobian-Vermutung: Ein falsches Analogon}
\label{abschn:real}
%% =============================================================================

Das natürliche reelle Analogon lautet: Ist $F \colon \mathbb{R}^n \to \mathbb{R}^n$
polynomial mit $\det(JF)(x) \neq 0$ für alle $x$, so ist $F$ bijektiv?

\begin{satz}[Pinchuk, 1994]\label{satz:pinchuk}
Die reelle Jacobian-Vermutung ist für $n = 2$ \emph{falsch}. Es existiert eine
polynomiale Abbildung $F \colon \mathbb{R}^2 \to \mathbb{R}^2$ mit:
\begin{enumerate}[label=(\roman*)]
  \item $\det(JF)(x) > 0$ für alle $x \in \mathbb{R}^2$,
  \item $F$ ist nicht surjektiv.
\end{enumerate}
\end{satz}

\begin{proof}[Beweisidee]
Pinchuk~\cite{Pin94} konstruiert explizit ein Polynom $Q(x,y)$ ohne reelle Nullstellen,
aber mit einer komplexen Kurve $\{Q = 0\} \subset \mathbb{C}^2$. Die Abbildung
$F = (P/Q, \ldots)$ (nach Beseitigung des Nenners über rationale Parametersubstitution
zu einem Polynom gemacht) hat positives Jacobi-Determinante überall, aber verfehlt
gewisse reelle Punkte aufgrund des algebraisch-geometrischen Verhaltens der Kurve.
\end{proof}

\begin{korollar}
Die komplexe Jacobian-Vermutung kann nicht durch rein reelle Methoden bewiesen werden.
\end{korollar}

%% =============================================================================
\section{Teilresultate für \texorpdfstring{$n=2$}{n=2}}
\label{abschn:n2}
%% =============================================================================

\begin{satz}[Wang, 1980]\label{satz:wang}
Die Jacobian-Vermutung gilt für alle polynomialen Abbildungen
$F \colon \mathbb{C}^n \to \mathbb{C}^n$ vom Grad $\leq 2$.
\end{satz}

\begin{satz}[Abhyankar--Moh--Suzuki]\label{satz:AMS}
Sind $P, Q \in \mathbb{C}[x,y]$ mit $\mathbb{C}[P,Q] = \mathbb{C}[x,y]$, dann gilt
$\deg P \mid \deg Q$ oder $\deg Q \mid \deg P$.
\end{satz}

Dieses Teilbarkeitsresultat für Grade erklärt, warum Jacobian-Abbildungen in
niedrigen Graden invertierbar sein müssen.

%% =============================================================================
\section{Äquivalenz mit der Dixmier-Vermutung}
\label{abschn:dixmier}
%% =============================================================================

\begin{definition_de}
Die \emph{Weyl-Algebra} $A_n(\mathbb{C})$ ist die $\mathbb{C}$-Algebra, erzeugt von
$x_1,\ldots,x_n, \partial_1, \ldots, \partial_n$ mit den Relationen
$[\partial_i, x_j] = \delta_{ij}$, $[x_i, x_j] = [\partial_i, \partial_j] = 0$.
Sie ist die Algebra der polynomialen Differentialoperatoren auf $\mathbb{C}^n$.
\end{definition_de}

\begin{vermutung}[Dixmier, 1968]\label{verm:dixmier}
Jeder Endomorphismus von $A_n(\mathbb{C})$ ist ein Automorphismus.
\end{vermutung}

\begin{satz}[Tsuchimoto 2005; Belov-Kanel--Kontsevich 2007]\label{satz:JV=DV}
Die Jacobian-Vermutung (Vermutung~\ref{verm:JV}) und die Dixmier-Vermutung
(Vermutung~\ref{verm:dixmier}) sind stabil äquivalent: $\mathrm{JV}_n$ impliziert
$\mathrm{DV}_n$, und $\mathrm{DV}_{2n}$ impliziert $\mathrm{JV}_n$.
\end{satz}

\begin{proof}[Beweisskizze ($\mathrm{JV} \Rightarrow \mathrm{DV}$)]
Sei $\varphi$ ein Endomorphismus von $A_n(\mathbb{C})$. Setze $P_i = \varphi(x_i)$
und $Q_i = \varphi(\partial_i)$. Die Kommutationsrelationen erzwingen
$\det(JF) = 1$ modulo Termen niedrigerer Ordnung (über den graduierten Ring).
Falls $\mathrm{JV}$ gilt, ist $F$ invertierbar, und das inverse Polynom liefert
den inversen Endomorphismus.
\end{proof}

\begin{bemerkung}
Die „stabile" Äquivalenz bedeutet: $\mathrm{JV}$ gilt für alle $n$ genau dann,
wenn $\mathrm{DV}$ für alle $n$ gilt. Dies ist eine tiefe Verbindung zwischen
affiner algebraischer Geometrie und nicht-kommutativer Algebra.
\end{bemerkung}

%% =============================================================================
\section{Weitere algebraische Äquivalenzen}
\label{abschn:equiv}
%% =============================================================================

Über die Dixmier-Vermutung hinaus ist die Jacobian-Vermutung äquivalent zu weiteren
offenen Problemen und unterstreicht damit ihre zentrale Rolle in der Mathematik.

\begin{vermutung}[Poisson-Vermutung, äquivalent]\label{verm:poisson_equiv}
Die Jacobian-Vermutung ist äquivalent dazu, dass jeder Endomorphismus der
Poisson-Algebra $\mathbb{C}[x_1,\ldots,x_n, y_1,\ldots,y_n]$ (mit Poisson-Klammer
$\{x_i, y_j\} = \delta_{ij}$), der die Poisson-Klammer erhält, ein Automorphismus ist.
\end{vermutung}

\begin{vermutung}[Stabile Zahmheit, äquivalent]\label{verm:stabile_zahmheit}
Ein Endomorphismus von $\mathbb{C}^n$, der die Jacobian-Bedingung erfüllt, ist genau
dann ein Automorphismus, wenn er \emph{stabil zahm} ist: nach Hinzufügen hinreichend
vieler Variablen wird er zu einer Komposition elementarer Automorphismen.
\end{vermutung}

%% =============================================================================
\section{Warum Standardmethoden versagen}
%% =============================================================================

\begin{bemerkung}[Versagen topologischer Argumente]
Über $\mathbb{C}$ ist eine eigentliche lokale Homöomorphie eine Überlagerungsabbildung.
Da $\mathbb{C}^n$ einfach zusammenhängend ist, wäre $F$ bijektiv. Das Problem: Es
ist \emph{nicht bekannt}, ob Jacobian-Abbildungen eigentlich sind, und dies ist
äquivalent zur Vermutung selbst.
\end{bemerkung}

\begin{bemerkung}[Versagen von Gradargumenten]
Das formale Inverse $G = \sum_{k \geq 0} (-H)^{\circ k}$ (iterierte Komposition)
ist eine formale Potenzreihe. Sie bricht genau dann zu einem Polynom ab, wenn
$(JH)^k = 0$ für alle hinreichend großen $k$. Die BCW-Reduktion auf $(JH)^2 = 0$
sucht genau diese strenge Nilpotenz zu erzwingen.
\end{bemerkung}

%% =============================================================================
\section{Verallgemeinerungen und verwandte Probleme}
%% =============================================================================

\begin{vermutung}[Mathieu-Vermutung, verwandt]\label{verm:mathieu}
Für Lie-Gruppen $G$ mit $G$-invariantem Differentialoperator $\Delta$: Falls
$f \in C^\infty(G)$ und $\Delta^k f = 0$ für alle großen $k$, muss $f$ eine
endliche Summe von Exponentialfunktionen sein? Dies impliziert via
Mathieus Substitutionstheorem~\cite{Math97} die Jacobian-Vermutung.
\end{vermutung}

%% Poisson-Vermutung bereits in Abschnitt~\ref{abschn:equiv} als äquivalente Formulierung
%% aufgeführt (Vermutung~\ref{verm:poisson_equiv}); Duplikat entfernt.

%% =============================================================================
\section{Offene Probleme und aktueller Forschungsstand}
\label{abschn:open}
%% =============================================================================

Die Jacobian-Vermutung bleibt eines der bekanntesten offenen Probleme der Mathematik.
Wir listen die primären offenen Fragen auf:

\begin{enumerate}[label=(\arabic*)]
  \item \textbf{Jacobian-Vermutung für $n = 2$}: Offen seit 1939. Selbst für
        Polynomabbildungen vom Grad $17$ über $\mathbb{C}$ ist die Vermutung nicht in
        voller Allgemeinheit bekannt.

  \item \textbf{Injektivität durch Jacobian}: Impliziert $\det(JF) = c \neq 0$
        bereits die Injektivität von $F$? Injektivität zusammen mit lokaler
        Diffeomorphie und Surjektivität würde die Vermutung ergeben; jede einzelne
        Implikation ist offen.

  \item \textbf{Reduktion auf Normalformen}: Kann jede Jacobian-Abbildung durch
        Komposition mit bekannten Automorphismen auf eine berechenbare Normalform
        gebracht werden, deren Invertierbarkeit prüfbar ist?

  \item \textbf{Algorithmik}: Ist es entscheidbar, ob eine gegebene polynomiale
        Abbildung für einen vorgegebenen Gradschranke $d$ ein polynomialer
        Automorphismus ist?

  \item \textbf{Positive Charakteristik}: Über $\mathbb{F}_p$ versagt die
        Jacobian-Vermutung (wie in Abschnitt~\ref{abschn:n1} für $n=1$ bemerkt).
        Was ist das korrekte Analogon?
\end{enumerate}

\begin{vermutung}[Verallgemeinerte Bass--Connell--Wright-Vermutung]\label{verm:genBCW}
Die Jacobian-Vermutung für Abbildungen vom Grad $n$ in Dimension $d$ ist äquivalent
zur Jacobian-Vermutung für Abbildungen vom Grad $3$ in Dimension
$d \cdot \binom{n+d-1}{d}$.
\end{vermutung}

%% =============================================================================
\section{Zusammenfassung}
%% =============================================================================

Die Jacobian-Vermutung steht als Denkmal für die Tiefe globaler Probleme in der
algebraischen Geometrie. Ihre elementare Formulierung verbirgt eine Tiefe, die den
besten Methoden der kommutativen Algebra, der komplexen Analysis und der algebraischen
Topologie über mehr als acht Jahrzehnte widerstanden hat.

Die wichtigsten bekannten Resultate:
\begin{itemize}
  \item $n = 1$: Trivial wahr (lineare Abbildungen).
  \item Grad $\leq 2$: Wahr (Wangs Theorem).
  \item Über $\mathbb{R}$, $n \geq 2$: \emph{Falsch} (Pinchuk 1994).
  \item Äquivalent zur Dixmier-Vermutung (Tsuchimoto, Belov-Kanel--Kontsevich).
  \item Reduziert auf kubisch-linearen nilpotenten Fall (BCW 1982).
\end{itemize}

Das Problem ist \textbf{offen für $n \geq 2$ über $\mathbb{C}$}.

\begin{thebibliography}{99}

\bibitem{Keller1939}
O.~H.\ Keller,
\textit{Ganze Cremona-Transformationen},
Monatshefte Math.\ Phys.\ \textbf{47} (1939), 299--306.

\bibitem{BCW82}
H.\ Bass, E.~H.\ Connell, D.\ Wright,
\textit{The Jacobian conjecture: reduction of degree and formal expansion of the inverse},
Bull.\ Amer.\ Math.\ Soc.\ (N.S.)\ \textbf{7} (1982), Nr.\ 2, 287--330.

\bibitem{Yag80}
A.~W.\ Yagzhev,
\textit{Über Kellers Problem} [russ.],
Sibirsk.\ Mat.\ Zh.\ \textbf{21} (1980), Nr.\ 5, 141--150.

\bibitem{Pin94}
S.~I.\ Pinchuk,
\textit{A counterexample to the strong real Jacobian conjecture},
Math.\ Z.\ \textbf{217} (1994), Nr.\ 1, 1--4.

\bibitem{ShU04}
I.~P.\ Shestakov, U.~U.\ Umirbaev,
\textit{The tame and the wild automorphisms of polynomial rings in three variables},
J.\ Amer.\ Math.\ Soc.\ \textbf{17} (2004), Nr.\ 1, 197--227.

\bibitem{Dixmier1968}
J.\ Dixmier,
\textit{Sur les algèbres de Weyl},
Bull.\ Soc.\ Math.\ France \textbf{96} (1968), 209--242.

\bibitem{Tsu05}
Y.\ Tsuchimoto,
\textit{Endomorphisms of Weyl algebra and $p$-curvatures},
Osaka J.\ Math.\ \textbf{42} (2005), Nr.\ 2, 435--452.

\bibitem{BK07}
A.\ Belov-Kanel, M.\ Kontsevich,
\textit{The Jacobian conjecture is stably equivalent to the Dixmier conjecture},
Mosc.\ Math.\ J.\ \textbf{7} (2007), Nr.\ 2, 209--218.

\bibitem{Wang80}
S.~S.-S.\ Wang,
\textit{A Jacobian criterion for separability},
J.\ Algebra \textbf{65} (1980), Nr.\ 2, 453--494.

\bibitem{Math97}
O.\ Mathieu,
\textit{Some conjectures about invariant theory and their applications},
in \textit{Algèbre non commutative}, Sémin.\ Congr.\ \textbf{2} (1997), 263--279.

\bibitem{vanEssen}
A.\ van den Essen,
\textit{Polynomial Automorphisms and the Jacobian Conjecture},
Progress in Mathematics \textbf{190}, Birkhäuser, Basel, 2000.

\bibitem{Jung42}
H.~W.~E.\ Jung,
\textit{Über ganze birationale Transformationen der Ebene},
J.\ Reine Angew.\ Math.\ \textbf{184} (1942), 161--174.

\end{thebibliography}

\end{document}
